\documentclass{umk-matematika}

\begin{document}

\maketitlepage
    {Практическая работа №2}
    {Значения тригонометрических функций}
    {}

\section{Фундамент (1--4 балла)}

\textbf{Задача 1 (1 балл).} Вычислите: $\sin 30^\circ$

\textbf{Задача 2 (2 балла).} Вычислите: $\cos 60^\circ$

\textbf{Задача 3 (3 балла).} Вычислите: $\tg 45^\circ$

\textbf{Задача 4 (4 балла).} Вычислите: $\sin 90^\circ$

\section{Стены (5--7 баллов)}

\textbf{Задача 5 (5 баллов).} Вычислите: $\sin 120^\circ$

\textbf{Задача 6 (6 баллов).} Вычислите: $\cos 150^\circ$

\textbf{Задача 7 (7 баллов).} Вычислите: $\tg 135^\circ$

\section{Кровля (8--10 баллов)}

\textbf{Задача 8 (8 баллов).} Вычислите без калькулятора: $\sin 150^\circ + \cos 120^\circ$

\textbf{Задача 9 (9 баллов).} Определите знак выражения: $\sin 200^\circ \cdot \cos 100^\circ$

\textbf{Задача 10 (10 баллов).} Кровля имеет уклон 30$^\circ$. Длина ската (гипотенуза) = 10 м. Найдите высоту подъёма кровли.

\newpage
\section*{Ответы}

\begin{enumerate}
  \item $\dfrac{1}{2}$
  \item $\dfrac{1}{2}$
  \item 1.
  \item 1.
  \item $\dfrac{\sqrt{3}}{2}$
  \item $-\dfrac{\sqrt{3}}{2}$
  \item –1.
  \item 0.
  \item плюс (+).
  \item высота подъёма = 5 м.
\end{enumerate}

\end{document}